% 大学物理の問題集 PDF版 共通スタイル

\usepackage[haranoaji]{luatexja-preset}
\usepackage{amsmath,amssymb,mathtools}
\usepackage{tikz}
\usepackage[most]{tcolorbox}
\usepackage{xcolor}
\usepackage{geometry}
\usepackage{fancyhdr}
\usepackage{enumitem}
\usepackage{hyperref}
\usepackage{bookmark}

\geometry{
  a4paper,
  top=23mm,
  bottom=24mm,
  inner=23mm,
  outer=20mm,
  headheight=28pt,
  headsep=5mm
}

\definecolor{PhysicsNavy}{HTML}{151E3D}
\definecolor{PhysicsBlueGray}{HTML}{E9EDF4}
\definecolor{PhysicsRule}{HTML}{AAB3C3}
\definecolor{PhysicsText}{HTML}{20242C}

\color{PhysicsText}
\setlength{\parindent}{1em}
\setlength{\parskip}{0.25\baselineskip}
\linespread{1.12}

\hypersetup{
  colorlinks=true,
  linkcolor=PhysicsNavy,
  urlcolor=PhysicsNavy,
  citecolor=PhysicsNavy,
  bookmarksnumbered=true
}

% 章・節見出し（jlreq純正の設定を利用）
\ModifyHeading{chapter}{
  font={\Huge\bfseries\color{PhysicsNavy}},
  label_format={第\thechapter 章}
}
\ModifyHeading{section}{
  font={\Large\bfseries\color{PhysicsNavy}}
}

% ヘッダー・フッター
\pagestyle{fancy}
\fancyhf{}
\fancyhead[L]{\small\color{PhysicsNavy}\leftmark}
\fancyhead[R]{\small\color{PhysicsNavy}大学物理の問題集}
\fancyfoot[C]{\small\color{PhysicsNavy}\thepage}
\renewcommand{\headrulewidth}{0.4pt}
\renewcommand{\headrule}{\hbox to\headwidth{\color{PhysicsRule}\leaders\hrule height \headrulewidth\hfill}}
\renewcommand{\chaptermark}[1]{\markboth{第\thechapter 章\quad #1}{}}

% 問題ボックス
\newtcolorbox{problemBox}[2][]{
  enhanced,
  breakable,
  colback=white,
  colframe=PhysicsNavy,
  colbacktitle=PhysicsNavy,
  coltitle=white,
  title={#2},
  fonttitle=\bfseries,
  boxrule=0.7pt,
  arc=1.2mm,
  left=5mm,
  right=5mm,
  top=4mm,
  bottom=4mm,
  before skip=0.5\baselineskip,
  after skip=\baselineskip,
  attach boxed title to top left={xshift=0mm},
  boxed title style={
    sharp corners=south,
    boxrule=0pt,
    left=4mm,
    right=4mm,
    top=1.6mm,
    bottom=1.6mm
  },
  #1
}

\newcommand{\solutionheading}{%
  \section*{解説}%
  \addcontentsline{toc}{section}{解説}%
  \noindent\color{PhysicsNavy}\rule{\linewidth}{0.8pt}%
  \par\vspace{0.9\baselineskip}\color{PhysicsText}%
}

\newcommand{\solutionpart}[1]{%
  \par\vspace{1.1\baselineskip}%
  \noindent{\large\bfseries\color{PhysicsNavy}#1}%
  \par\vspace{0.35\baselineskip}\color{PhysicsText}%
}

\setlist[enumerate]{
  label=(\arabic*),
  leftmargin=3em,
  itemsep=0.75\baselineskip,
  topsep=0.6\baselineskip
}

\setcounter{tocdepth}{1}
